% !TEX program = xelatex
\documentclass[12pt,a4paper]{article}
\usepackage[UTF8]{ctex}
\usepackage[left=2.2cm,right=2.2cm,top=1.8cm,bottom=1.8cm]{geometry}
\usepackage[table]{xcolor}
\usepackage{array}
\usepackage{tabularx}
\usepackage{float}
\usepackage{fancyhdr}
\usepackage{hyperref}
\usepackage{listings}
\usepackage{enumitem}
\usepackage{graphicx}

\newcommand{\StudentName}{张琳婕 (Zhang Linjie)}
\newcommand{\StudentID}{24020036086}
\newcommand{\ReportDate}{September 29, 2026}
\newcommand{\CourseName}{Robot Vision}

\definecolor{MainGreen}{HTML}{2E6B4E}
\definecolor{LightGreen}{HTML}{EEF5F0}
\definecolor{LighterGreen}{HTML}{F4F8F5}
\definecolor{MainText}{HTML}{1F2933}
\definecolor{GrayText}{HTML}{68737D}
\definecolor{TableLine}{HTML}{D9E1E5}
\colorlet{labgreen}{MainGreen}
\colorlet{lightgreen}{LightGreen}
\color{MainText}

\setlength{\parindent}{2em}
\setlength{\parskip}{0.4em}
\linespread{1.3}
\newcommand{\reportsection}[2]{\vspace{1.1em}\noindent{\fontsize{15pt}{19pt}\selectfont\bfseries\color{MainText}#1\quad #2}\par\vspace{0.6em}}
\newcommand{\reportsubsection}[2]{\vspace{0.7em}\noindent{\fontsize{14pt}{18pt}\selectfont\bfseries\color{MainGreen}#1\quad #2}\par\vspace{0.3em}}

\pagestyle{fancy}
\fancyhf{}
\fancyhead[R]{\small\color{GrayText}LAB 01\quad Robot Vision}
\fancyfoot[C]{\small\color{GrayText}Linux, Git, C++ and CMake Report\quad\textperiodcentered\quad\thepage}
\renewcommand{\headrulewidth}{0pt}
\renewcommand{\footrulewidth}{0pt}
\setlength{\headheight}{14pt}

\newcolumntype{Y}{>{\centering\arraybackslash}X}
\renewcommand{\arraystretch}{1.45}
\arrayrulecolor{TableLine}
\lstset{
  basicstyle=\ttfamily\small,
  breaklines=true,
  columns=fullflexible,
  frame=single,
  rulecolor=\color{TableLine},
  backgroundcolor=\color{LighterGreen},
  xleftmargin=0.5em,
  xrightmargin=0.5em,
  aboveskip=0.5em,
  belowskip=0.5em,
  showstringspaces=false
}
\hypersetup{colorlinks=true,urlcolor=MainGreen,linkcolor=MainGreen}

\begin{document}

{\bfseries\color{MainGreen} LAB REPORT / ROBOT VISION\par}
\vspace{1.2em}
{\fontsize{16pt}{21pt}\selectfont\bfseries Linux, Git, C++ and CMake Fundamentals\par}
\vspace{0.5em}
{\fontsize{14pt}{18pt}\selectfont\color{GrayText} Command-line practice and a RandomVector implementation\par}
\vspace{1.8em}

\begin{tabularx}{\textwidth}{>{\columncolor{LightGreen}\bfseries\color{GrayText}}p{2.4cm}X>{\columncolor{LightGreen}\bfseries\color{GrayText}}p{2.4cm}X}
\hline
Name & \StudentName & Student ID & \StudentID \\
\hline
Date & \ReportDate & Course & \CourseName \\
\hline
\end{tabularx}

\reportsection{01}{Experiment Overview}
This experiment introduces the Linux command line, Git, C++ and CMake. The work was completed on Ubuntu 22.04. A public course repository was cloned, text-processing and output-redirection exercises were completed, and the supplied \texttt{RandomVector} class was implemented without using algorithms from the \texttt{<algorithm>} header.

All required programs compiled and ran successfully. The C++ program generated 20 reproducible random values, calculated their mean and extrema, and printed a five-bin histogram. A separate CMake demonstration verified a multi-file executable linked against a custom library.

\begin{table}[H]
\centering
\rowcolors{2}{lightgreen}{white}
\begin{tabular}{p{4.0cm}p{3.2cm}p{4.0cm}}
\rowcolor{labgreen}\color{white}\bfseries Task & \color{white}\bfseries Result & \color{white}\bfseries Conclusion \\
Shell exercises & Completed & Output files verified \\
RandomVector & Compiled and executed & Statistics correct \\
CMake demo & Compiled and executed & Library linked successfully \\
\end{tabular}
\end{table}

\reportsection{02}{Experimental Environment}
\begin{table}[H]
\centering
\rowcolors{2}{lightgreen}{white}
\begin{tabular}{p{5cm}p{8cm}}
\rowcolor{labgreen}\color{white}\bfseries Item & \color{white}\bfseries Configuration \\
Operating System & Ubuntu 22.04 LTS \\
C++ Compiler & g++ 11.4.0 \\
CMake & 3.22.1 \\
Git & 2.34.1 \\
Build Standard & C++11 with \texttt{-Wall -pedantic} \\
\end{tabular}
\end{table}

Ubuntu 22.04 is sufficient for the Shell, Git, C++ and CMake content in Lab 1. If later laboratories require ROS 1 Noetic, a separate Ubuntu 20.04 environment should be used to match the course dependencies.

\clearpage
\reportsection{03}{Git and Shell Exercises}
\reportsubsection{3.1}{Course Repository}
The course repository was cloned from \url{https://github.com/ouc-vlab-course/vnav-codes}. The first attempt failed because of an HTTP/2 framing error. Retrying with HTTP/1.1 and a shallow clone completed successfully.

\begin{lstlisting}
git -c http.version=HTTP/1.1 clone --depth 1 \
  https://github.com/ouc-vlab-course/vnav-codes.git
\end{lstlisting}

\reportsubsection{3.2}{Text Statistics}
The file \texttt{dante.txt} was analysed with \texttt{wc} and \texttt{grep}. The first result line used \texttt{>} to create the output file, while subsequent lines used \texttt{>>} to append results.

\begin{table}[H]
\centering
\rowcolors{2}{lightgreen}{white}
\begin{tabular}{p{6cm}p{6cm}}
\rowcolor{labgreen}\color{white}\bfseries Metric & \color{white}\bfseries Result \\
Total lines & 19,567 \\
Total words & 97,676 \\
Nonblank lines & 14,338 \\
\end{tabular}
\end{table}

\begin{lstlisting}
echo "Lines: $(wc -l < dante.txt)" > exercise1.txt
echo "Words: $(wc -w < dante.txt)" >> exercise1.txt
echo "Nonblank lines: $(grep -cv '^[[:space:]]*$' dante.txt)" \
  >> exercise1.txt
\end{lstlisting}

\reportsubsection{3.3}{Append Redirection}
The \texttt{fortune} command was run five times. Every result was appended to \texttt{fortunes.txt}; the final file contained five lines. This demonstrates that \texttt{>>} preserves existing data, unlike \texttt{>}, which overwrites it.

\reportsection{04}{C++ Fundamentals}
\begin{table}[H]
\centering
\rowcolors{2}{lightgreen}{white}
\begin{tabular}{p{6.5cm}p{5.5cm}}
\rowcolor{labgreen}\color{white}\bfseries Question & \color{white}\bfseries Result \\
Prefix and postfix increment & $i=2$, $j=1$ \\
Conditional operator & Prints \texttt{b} \\
Reference assignment & Prints \texttt{10 10} \\
Pointer multiplication & Prints \texttt{1764} \\
Pointer and array indexing & Prints \texttt{0} \\
Reference parameter reset & Prints \texttt{0} \\
Conversion of 3.14 to \texttt{int} & Stores \texttt{3} \\
Conversion of $-1$ to 8-bit \texttt{unsigned char} & Stores \texttt{255} \\
Integer used as a condition & Final value is \texttt{0} \\
\end{tabular}
\end{table}

\clearpage
\reportsection{05}{RandomVector Implementation}
\reportsubsection{5.1}{Methods}
The constructor fills a vector with random values between zero and \texttt{max\_val}. The \texttt{mean()}, \texttt{min()} and \texttt{max()} methods use explicit loops. The histogram divides the interval from the minimum to the maximum into equal-width bins. A value equal to the maximum is assigned to the final bin to prevent an out-of-range index.

The implementation also checks invalid vector sizes, negative maximum values and non-positive bin counts. No function from \texttt{<algorithm>} is used.

\begin{lstlisting}
g++ -std=c++11 -Wall -pedantic \
  -o random_vector main.cpp random_vector.cpp
./random_vector
\end{lstlisting}

\reportsubsection{5.2}{Experimental Results}
The program uses the fixed seed 314159, making the sequence reproducible on the same C library implementation.

\begin{table}[H]
\centering
\rowcolors{2}{lightgreen}{white}
\begin{tabular}{p{6cm}p{6cm}}
\rowcolor{labgreen}\color{white}\bfseries Metric & \color{white}\bfseries Result \\
Number of samples & 20 \\
Mean & 0.558040 \\
Minimum & 0.0278694 \\
Maximum & 0.982655 \\
Five-bin counts & 5, 1, 4, 3, 7 \\
Remaining TODO markers & 0 \\
\end{tabular}
\end{table}

\begin{lstlisting}
Histogram:
                ***
                ***
***             ***
***     ***     ***
***     *** *** ***
***     *** *** ***
*** *** *** *** ***
\end{lstlisting}

\reportsection{06}{CMake Multi-file Build}
The demonstration project contains \texttt{main.cpp}, \texttt{my\_lib.cpp}, \texttt{my\_lib.hpp} and \texttt{CMakeLists.txt}. CMake creates the \texttt{mylib} library, builds the \texttt{main} executable and links the executable to the library.

\begin{lstlisting}
mkdir build && cd build
cmake ..
cmake --build .
./main

Hello world!
In library
\end{lstlisting}

The output confirms that the executable called the function implemented in the separate library source file.

\clearpage
\reportsection{07}{Analysis and Discussion}
The main setup difficulty was the interrupted GitHub connection. Forcing HTTP/1.1 avoided the HTTP/2 framing error without changing the repository. This demonstrates that a failed clone does not necessarily indicate an incorrect Git installation.

The most important programming boundary case occurred in histogram binning. The direct expression for a maximum-valued sample can produce an index equal to the number of bins. Since valid indices end at \texttt{bins - 1}, the implementation explicitly maps this value to the final bin. Checks for invalid constructor arguments and histogram parameters further improve robustness.

The measured Shell results are deterministic for the downloaded text. The RandomVector results are reproducible within the current environment because the program uses a fixed seed. Exact pseudo-random sequences can differ between C library implementations, so the algorithm and statistical properties are more important than matching every generated value on another platform.

\reportsection{08}{GitHub Repository and Commit History}
The completed exercises and source code are maintained in the personal public repository
\url{https://github.com/Eva122205/vnav-personal}. The Lab 1 materials are located in the
\texttt{lab1/} directory. Work was recorded through multiple Git commits instead of a single
final upload. Figure~\ref{fig:github-commits} shows the repository commit history at the time
of submission.

\begin{figure}[H]
  \centering
  \includegraphics[width=\textwidth]{assets/github_commits.png}
  \caption{GitHub commit history for Lab 1, including commit identifiers \texttt{8cf1cee} and \texttt{a6b4723}.}
  \label{fig:github-commits}
\end{figure}

\reportsection{09}{Conclusion}
This experiment successfully completed the introductory Linux, Git, C++ and CMake tasks. Shell commands produced the required text statistics and append-only output file. The C++ questions reinforced expression evaluation, references, pointers and numeric conversion. RandomVector compiled without warnings, calculated the expected statistics and produced a five-bin histogram. Finally, the CMake demonstration generated and built a multi-file project linked to a custom library.

The results show that Ubuntu 22.04 can complete every Lab 1 task. The complete work is stored under \texttt{\~/vnav-personal/lab1} for later submission.

\reportsection{10}{Submission Record}
\begin{tabular}{p{4.5cm}p{8.5cm}}
\rowcolor{labgreen}\color{white}\bfseries Item & \color{white}\bfseries Record \\
\rowcolor{lightgreen}Personal Repository & \url{https://github.com/Eva122205/vnav-personal} \\
\rowcolor{lightgreen}Course Repository & \url{https://github.com/ouc-vlab-course/vnav-codes} \\
Commit Records & \texttt{8cf1cee}, \texttt{a6b4723} \\
Shell Results & \texttt{exercise1.txt}, \texttt{fortunes.txt} \\
Written Answers & \texttt{cpp\_answers.txt} \\
C++ Program & \texttt{RandomVector/} \\
CMake Program & \texttt{lab1\_demo/} \\
\end{tabular}

\end{document}
